\begin{abstract}
\noindent
Constraint-based models such as flux balance analysis (FBA) predict cellular
metabolic states by solving linear optimization problems. When
nutrients or enzyme capacities vary continuously, optimal fluxes do
not change smoothly but follow piecewise-linear trajectories, changing slope
abruptly at bottlenecks. Here we show that the true
``curvature'' of this response is not an ordinary function but a
discrete measure concentrated on the boundary interfaces where active
constraints switch: the fundamental carrier of a network's rerouting
geometry. Path-dependent memory (holonomy) scales linearly with
perturbation size (slope $1.00$; smooth systems quadratic), and
$93.4$--$100.0\%$ of second-order adjustment concentrates at discrete
bottleneck transitions. A bridge connects these switches to smooth
differential geometry: under grid refinement they converge weakly to
smooth curvature densities, dual-cell reconstructions strongly in
$L^1$, but total variation fails to converge generically (grid
anisotropy), defining the resolution window ($h \ll \sigma \ll
L_{\mathrm{var}}$).

Integrating this measure along physiological paths yields a
parameter-free gene sensitivity metric, $\kmu$: each enzyme's
rerouting burden. In carbon-starved \emph{E.~coli}, $\kmu$ predicts
transcriptional log-fold changes across $424$ genes ($r = +0.395$,
$p = 2.6 \times 10^{-17}$; partial $r = +0.269$), robust across
tie-breaking rules and stress axes. The association vanishes at the
protein layer ($r = -0.083$ across $366$ genes in matched quantitative
proteomics): cells transcribe standby capacity for rerouting,
buffering protein synthesis. Double-knockout epistasis mirrors
active-set boundary overlaps ($\rho_S = 0.865$), and cyclic
perturbations leave lasting metabolic memory ($66\%$ non-reverting).
Our framework unites linear programming, discrete geometry, and
transcriptional regulation into an accessible, predictive foundation
for metabolic systems biology.
\end{abstract}
